% English translation, made from our own transcription (original.tex), not from any existing
% translation. All curve labels, the fraction, and \origpage markers are preserved unchanged;
% only the prose is translated. Period terms kept: "quantity of motion", "living force" (vis
% viva). As in the transcription, Gerhardt's editorial footnotes and the Huygens "Beilage" are
% not part of this work and are omitted.
\documentclass{article}
\usepackage{readmasters}
\begin{document}

\origpage{234}
\section*{III. On the isochronous curve, on which a heavy body descends without acceleration, and on the controversy with the Abbé De Conti}

When, in these Acta published in March 1686, I had brought out a demonstration against the
\emph{Cartesians}, in which the true estimation of forces is set forth and it is shown that it is
not the quantity of motion, but potentia, differing from the quantity of motion, that is
conserved, a
certain learned man in France, the Abbé \emph{De Conti}, replied on behalf of the Cartesians, but
— as afterwards appeared — without having sufficiently grasped the force of my argument. For he
believed that certain other received principles were being attacked by me, which he lists in the
Nouvelles de la République des Lettres for June 1687, p.~579, and he denies (p.~519 ff.) that he
recognises the contradiction which I seem to myself to find in them: although it never came into
my mind to doubt those principles, as I reminded him in the Nouvelles de la République des Lettres
for September 1687. In order to evade my objection, he had betaken himself to a side-issue of
time, which — in the way in which the state of the controversy had been conceived by me — is
plainly accidental. For, the height remaining the same, the same force is acquired or expended by
heavy bodies whatever the time allowed, which is increased or diminished according as the
inclination of the descent is greater or less. On that occasion, so that it might appear more
clearly that time — and hence the distinction between isochronous and non-isochronous powers — is
here nothing to the purpose, and that our dispute might yield some increase of knowledge, I
proposed to him, in the said Nouvelles for September 1687, the following problem, solved by me in
the course of writing and, as it seems, not inelegant: ``To find an isochronous curve on which a
heavy body descends uniformly — that is, in equal times draws equally near to the horizon — and so
is carried downward without acceleration and always with equal velocity.'' But the Abbé
\emph{De Conti} made no further reply, whether because he was unwilling to touch the problem, or
because, having at last understood my meaning, he judged himself satisfied. But in his stead the
most celebrated \emph{Christiaan Huygens} judged this problem worthy of his effort; his solution,
wholly consonant with mine, appears in the Nouvelles de la République des Lettres for October
1687, though with the demonstration suppressed and the explanation
\origpage{235}
of the distinction between the various lines of the same kind, as he says, which he notes as
satisfying it. These things, therefore, I wished to supply in this place, and would have done so
sooner, had I not been awaiting something here from the Abbé's industry.

\emph{Problem.} To find a plane curve on which a heavy body descends without acceleration.

\emph{Solution.} Let there be (fig.~116)\ednote{The plate containing fig.~116 could not be
located in the scans available for this work (the Gerhardt edition scan has no plates section);
it is not reproduced here.} any semicubical paraboloid curve $\beta Ne$ (namely, one
in which the solid under the square of the base $NM$ and the parameter $aP$ is equal to the cube
of the height $\beta M$), so placed that the tangent $\beta M$ at the vertex $\beta$ is
perpendicular to the horizon. If at any point $N$ whatever of this curve a heavy body be placed,
endowed with that further speed of descent which it could have acquired by descending from the
horizontal $Aa$ — whose elevation $a\beta$ above the vertex $\beta$ is $\frac{4}{9}$ of the
parameter of the curve — then the same body will descend further uniformly along the curve $Ne$,
however far continued, as was desired.

\emph{Demonstration.} Let the straight line $NT$ touch the curve $\beta Ne$ at $N$ and meet
$\beta M$ at $T$. Then certainly (from the known property of the tangents of this curve) $TM$ will
be to $NM$ in the subduplicate ratio of $a\beta$ to $\beta M$. Therefore $TM$ to $TN$ will be in
the subduplicate ratio of $a\beta$ to $a\beta + \beta M$, that is, to $aM$. Now the ratio of $TM$
to $TN$ is the same as that of the velocity of descending along the curve (that is, of drawing
nearer to the horizon along the curve) which the heavy body has when placed at $N$, to the
velocity with which the same body would descend from $N$ onward, not along the curve, but freely,
if it could (as is clear from the nature of inclined motion). But this free velocity is further
to a certain constant one in the subduplicate ratio of $aM$ to $a\beta$; for (as is clear from the
motion of heavy bodies) free velocities are in the subduplicate ratio of the heights
\origpage{236}
(from which, by descending, they are sought), namely $aM$: therefore the velocity of descending
along the curve $Ne$, which the body has when placed at any point $N$ of the curve, is to the
constant velocity in the ratio compounded of the subduplicate ratio of $a\beta$ to $aM$ and of
$aM$ to $a\beta$ — which is a ratio of equality. That very velocity of descending along the curve
is therefore constant, that is, everywhere the same on the curve $Ne$. Which was to be shown.

\emph{Corollaries:} 1) A heavy body having a speed as though fallen from some height $Aa$ can
descend from the same point $N$ along infinitely many isochronous curves, but of the same species
— differing only in the magnitude of the parameter — such as $Ne$, $N(e)$, $NE$, which are all
semicubical paraboloids, and therefore similar to one another. Indeed any paraboloid of this kind
serves here, provided it be so placed that $a\beta$ (or $(a)(\beta)$), the distance of the vertex
from the horizontal $a(a)$ from which the body began to descend, is $\frac{4}{9}$ of the parameter
of the curve $\beta e$ (or $(\beta)(e)$): nor does it matter whether the body, about to descend
isochronously along the curve $N(e)$, has reached $N$ from $a(a)$ by the path $(a)(\beta)N$, or by
some other, or has acquired the same speed and direction without any descent at all, from some
other cause. Yet of those infinitely many isochronous curves, along which the body can descend
from $N$ onward without acceleration, that one affords it the swiftest descent whose vertex is the
point $N$ itself — such as $NE$, which the straight line $AN$, perpendicular to the horizon,
touches.

2) The time of descent along the straight line $a\beta$ is to the time of descent along the curve
$\beta N$ as half the height $\beta M$ is to $a\beta$ itself: and accordingly, if $\beta M$ be
double $a\beta$, the times of the descents along $a\beta$ and along $\beta N$ will be equal. The
reason for this is manifest: for the times of uniform descent are to one another as the heights;
and, from what has been demonstrated by \emph{Galileo}, the time in which a moving body traverses
the height $a\beta$ by accelerated motion is double that in which it traverses an equal height
$\beta M$ (as happens here, though along the curve $\beta N$) by uniform motion having a speed
equal to the last one acquired by acceleration at $\beta$.

I confess, however, that I proposed this problem not to the foremost geometers, who are versed in
a certain deeper Analysis, but rather to those who side with \emph{that learned Gaul}, whom my
complaint about most present-day Cartesians (paraphrasts of the Master, rather than his rivals)
seemed to have somewhat offended. For such men attribute too much both to the dogmas received
among the Cartesians and to the analysis widely current among them, so much so that they suppose
themselves able, by its aid, to accomplish anything whatever in Mathematics
\origpage{237}
(provided only they are willing to take on the labour of calculation), not without detriment to
the sciences, which are now cultivated too negligently out of a false confidence in what has been
discovered. To these I had wished to offer, in this problem, matter for exercising their Analysis
— a problem that needs not lengthy calculation, but art.

If, however, anyone should complain that the solution has already been snatched from him, he may
seek another isochrone neighbouring on this one, in which the heavy body recedes uniformly not, as
hitherto, from the horizontal (or approaches it), but from a certain point. Whence the problem
will be this: \emph{to find a curve on which a descending heavy body recedes uniformly from a
given point, or approaches it.}

Such would be the curve $NQR$, if it were of such a nature that, from a given or fixed point $A$,
straight lines being drawn to the curve at will, such as $AN$, $AQ$, $AR$, the excess of $AR$
over $AQ$ were to the excess of $AQ$ over $AN$ as the time in which the arc $QR$ is descended is
to the time in which the arc $NQ$ is descended.

\end{document}
